\IfFileExists{stacks-project.cls}{%
\documentclass{stacks-project}
}{%
\documentclass{amsart}
}

% For dealing with references we use the comment environment
\usepackage{verbatim}
\newenvironment{reference}{\comment}{\endcomment}
%\newenvironment{reference}{}{}
\newenvironment{slogan}{\comment}{\endcomment}
\newenvironment{history}{\comment}{\endcomment}

% For commutative diagrams we use Xy-pic
\usepackage[all]{xy}

% We use 2cell for 2-commutative diagrams.
\xyoption{2cell}
\UseAllTwocells

% We use multicol for the list of chapters between chapters
\usepackage{multicol}

% This is generally recommended for better output
\usepackage{lmodern}
\usepackage[T1]{fontenc}

% For cross-file-references
\usepackage{xr-hyper}

% Package for hypertext links:
\usepackage{hyperref}

% For any local file, say "hello.tex" you want to link to please
% use \externaldocument[hello-]{hello}
\externaldocument[introduction-]{introduction}
\externaldocument[conventions-]{conventions}
\externaldocument[sets-]{sets}
\externaldocument[categories-]{categories}
\externaldocument[topology-]{topology}
\externaldocument[sheaves-]{sheaves}
\externaldocument[sites-]{sites}
\externaldocument[stacks-]{stacks}
\externaldocument[fields-]{fields}
\externaldocument[algebra-]{algebra}
\externaldocument[brauer-]{brauer}
\externaldocument[homology-]{homology}
\externaldocument[derived-]{derived}
\externaldocument[simplicial-]{simplicial}
\externaldocument[more-algebra-]{more-algebra}
\externaldocument[smoothing-]{smoothing}
\externaldocument[modules-]{modules}
\externaldocument[sites-modules-]{sites-modules}
\externaldocument[injectives-]{injectives}
\externaldocument[cohomology-]{cohomology}
\externaldocument[sites-cohomology-]{sites-cohomology}
\externaldocument[dga-]{dga}
\externaldocument[dpa-]{dpa}
\externaldocument[sdga-]{sdga}
\externaldocument[hypercovering-]{hypercovering}
\externaldocument[schemes-]{schemes}
\externaldocument[constructions-]{constructions}
\externaldocument[properties-]{properties}
\externaldocument[morphisms-]{morphisms}
\externaldocument[coherent-]{coherent}
\externaldocument[divisors-]{divisors}
\externaldocument[limits-]{limits}
\externaldocument[varieties-]{varieties}
\externaldocument[topologies-]{topologies}
\externaldocument[descent-]{descent}
\externaldocument[perfect-]{perfect}
\externaldocument[more-morphisms-]{more-morphisms}
\externaldocument[flat-]{flat}
\externaldocument[groupoids-]{groupoids}
\externaldocument[more-groupoids-]{more-groupoids}
\externaldocument[etale-]{etale}
\externaldocument[chow-]{chow}
\externaldocument[intersection-]{intersection}
\externaldocument[pic-]{pic}
\externaldocument[weil-]{weil}
\externaldocument[adequate-]{adequate}
\externaldocument[dualizing-]{dualizing}
\externaldocument[duality-]{duality}
\externaldocument[discriminant-]{discriminant}
\externaldocument[derham-]{derham}
\externaldocument[local-cohomology-]{local-cohomology}
\externaldocument[algebraization-]{algebraization}
\externaldocument[curves-]{curves}
\externaldocument[resolve-]{resolve}
\externaldocument[models-]{models}
\externaldocument[functors-]{functors}
\externaldocument[equiv-]{equiv}
\externaldocument[pione-]{pione}
\externaldocument[etale-cohomology-]{etale-cohomology}
\externaldocument[proetale-]{proetale}
\externaldocument[relative-cycles-]{relative-cycles}
\externaldocument[more-etale-]{more-etale}
\externaldocument[trace-]{trace}
\externaldocument[crystalline-]{crystalline}
\externaldocument[spaces-]{spaces}
\externaldocument[spaces-properties-]{spaces-properties}
\externaldocument[spaces-morphisms-]{spaces-morphisms}
\externaldocument[decent-spaces-]{decent-spaces}
\externaldocument[spaces-cohomology-]{spaces-cohomology}
\externaldocument[spaces-limits-]{spaces-limits}
\externaldocument[spaces-divisors-]{spaces-divisors}
\externaldocument[spaces-over-fields-]{spaces-over-fields}
\externaldocument[spaces-topologies-]{spaces-topologies}
\externaldocument[spaces-descent-]{spaces-descent}
\externaldocument[spaces-perfect-]{spaces-perfect}
\externaldocument[spaces-more-morphisms-]{spaces-more-morphisms}
\externaldocument[spaces-flat-]{spaces-flat}
\externaldocument[spaces-groupoids-]{spaces-groupoids}
\externaldocument[spaces-more-groupoids-]{spaces-more-groupoids}
\externaldocument[bootstrap-]{bootstrap}
\externaldocument[spaces-pushouts-]{spaces-pushouts}
\externaldocument[spaces-chow-]{spaces-chow}
\externaldocument[groupoids-quotients-]{groupoids-quotients}
\externaldocument[spaces-more-cohomology-]{spaces-more-cohomology}
\externaldocument[spaces-simplicial-]{spaces-simplicial}
\externaldocument[spaces-duality-]{spaces-duality}
\externaldocument[formal-spaces-]{formal-spaces}
\externaldocument[restricted-]{restricted}
\externaldocument[spaces-resolve-]{spaces-resolve}
\externaldocument[formal-defos-]{formal-defos}
\externaldocument[defos-]{defos}
\externaldocument[cotangent-]{cotangent}
\externaldocument[examples-defos-]{examples-defos}
\externaldocument[algebraic-]{algebraic}
\externaldocument[examples-stacks-]{examples-stacks}
\externaldocument[stacks-sheaves-]{stacks-sheaves}
\externaldocument[criteria-]{criteria}
\externaldocument[artin-]{artin}
\externaldocument[quot-]{quot}
\externaldocument[stacks-properties-]{stacks-properties}
\externaldocument[stacks-morphisms-]{stacks-morphisms}
\externaldocument[stacks-limits-]{stacks-limits}
\externaldocument[stacks-cohomology-]{stacks-cohomology}
\externaldocument[stacks-perfect-]{stacks-perfect}
\externaldocument[stacks-introduction-]{stacks-introduction}
\externaldocument[stacks-more-morphisms-]{stacks-more-morphisms}
\externaldocument[stacks-geometry-]{stacks-geometry}
\externaldocument[moduli-]{moduli}
\externaldocument[moduli-curves-]{moduli-curves}
\externaldocument[examples-]{examples}
\externaldocument[exercises-]{exercises}
\externaldocument[guide-]{guide}
\externaldocument[desirables-]{desirables}
\externaldocument[coding-]{coding}
\externaldocument[obsolete-]{obsolete}
\externaldocument[fdl-]{fdl}
\externaldocument[index-]{index}

% Theorem environments.
%
\theoremstyle{plain}
\newtheorem{theorem}[subsection]{Theorem}
\newtheorem{proposition}[subsection]{Proposition}
\newtheorem{lemma}[subsection]{Lemma}

\theoremstyle{definition}
\newtheorem{definition}[subsection]{Definition}
\newtheorem{example}[subsection]{Example}
\newtheorem{exercise}[subsection]{Exercise}
\newtheorem{situation}[subsection]{Situation}

\theoremstyle{remark}
\newtheorem{remark}[subsection]{Remark}
\newtheorem{remarks}[subsection]{Remarks}

\numberwithin{equation}{subsection}

% Macros
%
\def\lim{\mathop{\mathrm{lim}}\nolimits}
\def\colim{\mathop{\mathrm{colim}}\nolimits}
\def\Spec{\mathop{\mathrm{Spec}}}
\def\Hom{\mathop{\mathrm{Hom}}\nolimits}
\def\Ext{\mathop{\mathrm{Ext}}\nolimits}
\def\SheafHom{\mathop{\mathcal{H}\!\mathit{om}}\nolimits}
\def\SheafExt{\mathop{\mathcal{E}\!\mathit{xt}}\nolimits}
\def\Sch{\mathit{Sch}}
\def\Mor{\mathop{\mathrm{Mor}}\nolimits}
\def\Ob{\mathop{\mathrm{Ob}}\nolimits}
\def\Sh{\mathop{\mathit{Sh}}\nolimits}
\def\NL{\mathop{N\!L}\nolimits}
\def\CH{\mathop{\mathrm{CH}}\nolimits}
\def\proetale{{pro\text{-}\acute{e}tale}}
\def\etale{{\acute{e}tale}}
\def\QCoh{\mathit{QCoh}}
\def\Ker{\mathop{\mathrm{Ker}}}
\def\Im{\mathop{\mathrm{Im}}}
\def\Coker{\mathop{\mathrm{Coker}}}
\def\Coim{\mathop{\mathrm{Coim}}}

% Boxtimes
%
\DeclareMathSymbol{\boxtimes}{\mathbin}{AMSa}{"02}

%
% Macros for moduli stacks/spaces
%
\def\QCohstack{\mathcal{QC}\!\mathit{oh}}
\def\Cohstack{\mathcal{C}\!\mathit{oh}}
\def\Spacesstack{\mathcal{S}\!\mathit{paces}}
\def\Quotfunctor{\mathrm{Quot}}
\def\Hilbfunctor{\mathrm{Hilb}}
\def\Curvesstack{\mathcal{C}\!\mathit{urves}}
\def\Polarizedstack{\mathcal{P}\!\mathit{olarized}}
\def\Complexesstack{\mathcal{C}\!\mathit{omplexes}}
% \Pic is the operator that assigns to X its picard group, usage \Pic(X)
% \Picardstack_{X/B} denotes the Picard stack of X over B
% \Picardfunctor_{X/B} denotes the Picard functor of X over B
\def\Pic{\mathop{\mathrm{Pic}}\nolimits}
\def\Picardstack{\mathcal{P}\!\mathit{ic}}
\def\Picardfunctor{\mathrm{Pic}}
\def\Deformationcategory{\mathcal{D}\!\mathit{ef}}
\usepackage{cleveref}
\crefname{originaldefinition}{original definition}{original definitions}
\Crefname{originaldefinition}{Original definition}{Original definitions}
\makeatletter
\newlabel{definition-divided-powers-dga}{{09PJ}{source}{}{sourcepointer-1}{}}
\newlabel{definition-divided-powers-dga@cref}{{[originaldefinition][0][]09PJ}{[1][0][]source}{}{}{}}
\makeatother
\crefname{originalreference}{original reference}{original references}
\Crefname{originalreference}{Original reference}{Original references}
\makeatletter
\newlabel{example-adjoining-even}{{09PI}{source}{}{sourcepointer-2}{}}
\newlabel{example-adjoining-even@cref}{{[originalreference][0][]09PI}{[1][0][]source}{}{}{}}
\makeatother
\crefname{originalreference}{original reference}{original references}
\Crefname{originalreference}{Original reference}{Original references}
\makeatletter
\newlabel{example-adjoining-odd}{{09PH}{source}{}{sourcepointer-3}{}}
\newlabel{example-adjoining-odd@cref}{{[originalreference][0][]09PH}{[1][0][]source}{}{}{}}
\makeatother
\crefname{originallemma}{original lemma}{original lemmas}
\Crefname{originallemma}{Original lemma}{Original lemmas}
\makeatletter
\newlabel{lemma-extend-differential}{{09PM}{source}{}{sourcepointer-4}{}}
\newlabel{lemma-extend-differential@cref}{{[originallemma][0][]09PM}{[1][0][]source}{}{}{}}
\makeatother
\title{Tate resolutions: factorization through a divided power DGA quasi-isomorphism}
\author{The Stacks Project Authors}
\date{stacks-project:dpa.tex:local-snapshot-20261012; GNU FDL 1.2 or later}
\begin{document}
\maketitle
\noindent\textit{Editorial note.} 'Moreover, it is clear that H\_m(A(m+1)) = 0 by construction' rests on the cohomology computations of the cited adjoining Examples (in-chapter template). Declared dependency: lemma-extend-differential is not carded (its source proof is 'a direct computation which is omitted'); recorded via external\_references, not re-proved.\par\medskip
\noindent\textit{Allowed standard prerequisites (editorial).}\begin{itemize}
\item Differential extension to a DP variable. lemma-extend-differential (same chapter, named; NOT carded -- build-gate exclusion). Extends d uniquely to A(m) after adjoining variables with prescribed boundaries; declared as an external reference, not re-proved (the source proof of that lemma is omitted).
\item Adjoining DP variables, odd and even. example-adjoining-odd, example-adjoining-even (same chapter, named examples). Supply the divided power polynomial algebras in each degree and the cohomology bookkeeping behind 'H\_m(A(m+1)) = 0 by construction'.
\end{itemize}\medskip
\ifdefined\hypertarget\hypertarget{sourcepointer-1}{}\fi
\noindent\textit{Source reference.} Original definition 09PJ: sources/stacks-dpa/pinned-original-source.tex, source lines 1300--1300. Navigation pointer only: the counted proof cites this statement by \textbackslash{}ref as a standing input; the deduction the counted proof has to make is carried out inside the counted span, and no part of that deduction is delegated to this pointer.\par\medskip
\ifdefined\hypertarget\hypertarget{sourcepointer-2}{}\fi
\noindent\textit{Source reference.} Original reference 09PI: sources/stacks-dpa/pinned-original-source.tex, source lines 1241--1241. Navigation pointer only: the counted proof cites this statement by \textbackslash{}ref as a standing input; the deduction the counted proof has to make is carried out inside the counted span, and no part of that deduction is delegated to this pointer.\par\medskip
\ifdefined\hypertarget\hypertarget{sourcepointer-3}{}\fi
\noindent\textit{Source reference.} Original reference 09PH: sources/stacks-dpa/pinned-original-source.tex, source lines 1221--1221. Navigation pointer only: the counted proof cites this statement by \textbackslash{}ref as a standing input; the deduction the counted proof has to make is carried out inside the counted span, and no part of that deduction is delegated to this pointer.\par\medskip
\ifdefined\hypertarget\hypertarget{sourcepointer-4}{}\fi
\noindent\textit{Source reference.} Original lemma 09PM: sources/stacks-dpa/pinned-original-source.tex, source lines 1398--1398. Navigation pointer only: the counted proof cites this statement by \textbackslash{}ref as a standing input; the deduction the counted proof has to make is carried out inside the counted span, and no part of that deduction is delegated to this pointer.\par\medskip
\begin{lemma}
\label{lemma-tate-resolution}
Let $R \to S$ be a homomorphism of commutative rings.
There exists a factorization
$$
R \to A \to S
$$
with the following properties:
\begin{enumerate}
\item $(A, \text{d}, \gamma)$ is as in
Definition \ref{definition-divided-powers-dga},
\item $A \to S$ is a quasi-isomorphism (if we endow $S$ with
the zero differential),
\item $A_0 = R[x_j: j\in J] \to S$ is any surjection of a polynomial
ring onto $S$, and
\item $A$ is a graded divided power polynomial algebra over $R$.
\end{enumerate}
The last condition means that $A$ is constructed out of $A_0$ by
successively adjoining a set of variables $T$ in each degree $> 0$ as in
Example \ref{example-adjoining-odd} or \ref{example-adjoining-even}.
Moreover, if $R$ is Noetherian and $R\to S$ is of finite type,
then $A$ can be taken to have only finitely many generators in
each degree.
\end{lemma}

\begin{proof}
We write out the construction for the case that $R$ is Noetherian
and $R\to S$ is of finite type. Without those assumptions, the proof
is the same, except that we have to use some set (possibly
infinite) of generators in each degree.

\medskip\noindent
Start of the construction: Let $A(0) = R[x_1, \ldots, x_n]$ be
a (usual) polynomial ring and let $A(0) \to S$ be a surjection.
As grading we take $A(0)_0 = A(0)$ and $A(0)_d = 0$ for $d \not = 0$.
Thus $\text{d} = 0$ and $\gamma_n$, $n > 0$, is zero as well.

\medskip\noindent
Choose generators $f_1, \ldots, f_m \in R[x_1, \ldots, x_n]$
for the kernel of the given map $A(0) = R[x_1, \ldots, x_n] \to S$.
We apply Example \ref{example-adjoining-odd} $m$ times to get
$$
A(1) = A(0)\langle T_1, \ldots, T_m\rangle
$$
with $\deg(T_i) = 1$ as a graded divided power polynomial algebra.
We set $\text{d}(T_i) = f_i$. Since $A(1)$ is a divided power polynomial
algebra over $A(0)$ and since $\text{d}(f_i) = 0$
this extends uniquely to a differential on $A(1)$ by
Lemma \ref{lemma-extend-differential}.

\medskip\noindent
Induction hypothesis: Assume we are given factorizations
$$
R \to A(0) \to A(1) \to \ldots \to A(m) \to S
$$
where $A(0)$ and $A(1)$ are as above and each $R \to A(m') \to S$
for $2 \leq m' \leq m$ satisfies properties (1) and (4)
of the statement of the lemma and (2) replaced by the condition that
$H_i(A(m')) \to H_i(S)$ is an isomorphism for
$m' > i \geq 0$. The base case is $m = 1$.

\medskip\noindent
Induction step: Assume we have $R \to A(m) \to S$
as in the induction hypothesis. Consider the
group $H_m(A(m))$. This is a module over $H_0(A(m)) = S$.
In fact, it is a subquotient of $A(m)_m$ which is a finite
type module over $A(m)_0 = R[x_1, \ldots, x_n]$.
Thus we can pick finitely many elements
$$
e_1, \ldots, e_t \in \Ker(\text{d} : A(m)_m \to A(m)_{m - 1})
$$
which map to generators of this module. Applying
Example \ref{example-adjoining-odd} or
\ref{example-adjoining-even} $t$ times we get
$$
A(m + 1) = A(m)\langle T_1, \ldots, T_t\rangle
$$
with $\deg(T_i) = m + 1$ as a graded divided power algebra. We set
$\text{d}(T_i) = e_i$. Since $A(m+1)$ is a divided power polynomial
algebra over $A(m)$ and since $\text{d}(e_i) = 0$
this extends uniquely to a differential on $A(m + 1)$
compatible with the divided power structure.
Since we've added only material in degree $m + 1$ and higher we see
that $H_i(A(m + 1)) = H_i(A(m))$ for $i < m$. Moreover, it is
clear that $H_m(A(m + 1)) = 0$ by construction.

\medskip\noindent
To finish the proof we observe that we have shown there exists
a sequence of maps
$$
R \to A(0) \to A(1) \to \ldots \to A(m) \to A(m + 1) \to \ldots \to S
$$
and to finish the proof we set $A = \colim A(m)$.
\end{proof}

\end{document}
